\chapterimage{figs/chapterhead.png}
\chapter{GUI, Simulang, model execution and observation}
\label{chap:gui_simulang_execucao_observacao}

\section{Introduction}
\label{sec:cap5_introducao}

In the previous chapters, the reader got \textit{GenESyS} up and running, got to know the general logic of the graphical interface and started working with example models already saved in the project. From this point forward, it becomes possible to take an important step: move away from just a structural reading of the models and start observing them in execution, following their dynamics and relating the graphical representation to the textual representation of the system.

This chapter has exactly that purpose. It delves deeper into GUI usage in three complementary directions. First, it shows how the execution of the model should be followed by the user, not as an opaque operation, but as an observable process. Second, it introduces the \textit{Simulang} tab as a textual form of representation of the same model that has already been seen in the graphical interface. Third, it guides the observation of the first results, traces and signs of the system's behavior, so that the user begins to build a more informed reading of the simulation.

The importance of this stage is great. A really useful simulator is not only able to open models and run them, but to allow the user to understand what is happening. The \textit{GenESyS} GUI was designed to offer different perspectives of the same work: a graphical perspective, a textual perspective and a dynamic perspective, linked to execution. From here, the reader will begin to articulate these three layers.

\section{From static structure to model dynamics}
\label{sec:cap5_estrutura_estatica_dinamica}

When the user opens a model and looks at its graphical structure, he sees a static description of the system. The components, connections, supporting data and graphic elements indicate how the simulation was organized, but do not yet show, by themselves, the effective temporal behavior of the model. Running the simulation transforms this static structure into a dynamic process.

This point is central to understanding the simulator. A model is not just a well-organized diagram; he is a behavior machine. The function of the simulator is precisely to activate this machine, driving the advancement of simulated time, the processing of events, the circulation of entities and the updating of the system's relevant information.

For the user, the main consequence is this: execution should not be seen as just “pressing the start button”. It should be understood as the passage from the model description to its operational dynamics.

\subsection{What changes when the model starts running}
\label{subsec:cap5_o_que_muda_quando_executa}

When the simulation starts, the user starts working with a new type of reading. Before, the main question was:
\begin{quote}
How is this model organized?
\end{quote}

During execution, this question changes to:
\begin{quote}
How does this model behave over time?
\end{quote}

This means observing:
\begin{itemize}
\item whether entities are being created, transformed or removed;
\item if the flow actually goes through the elements that the model suggests;
\item whether resources, queues or supporting data appear to be used coherently;
\item if the global behavior of the simulation makes sense in relation to the previously read structure.
\end{itemize}

\section{Execution commands and their usage logic}
\label{sec:cap5_comandos_execucao_logica}

The \textit{GenESyS} GUI provides a set of commands directly linked to the simulation. Among them, the commands to start, pause, resume, step by step, stop and configure stand out. For the user, these commands should not be perceived just as operational buttons, but as different ways of relating to the dynamics of the model.

\subsection{Start the simulation}
\label{subsec:cap5_iniciar_simulacao}

The simulation start command is the transition point between the static model and observed behavior. Before activating it, the user must ensure that:
\begin{itemize}
\item the correct model is loaded;
\item the visual structure makes sense;
\item the system appears to be ready to run;
\item the objective of the observation is minimally clear.
\end{itemize}

This last observation is important. It is better to run the simulation with some reading hypothesis — for example, observing a specific flow or the effect of a certain property — than simply starting the execution without knowing what you want to track.

\subsection{Pause, resume and advance step by step}
\label{subsec:cap5_pausar_retormar_passo_a_passo}

One of the great advantages of a rich GUI is that it allows different rates of observation of the simulation. In some situations, the user just wants to check the overall behavior. Other times, you want to look at specific parts of the stream more closely. It is in this context that commands such as pause, resume and step by step become particularly valuable.

These commands transform the simulation into an inspection object. Instead of just letting the system run to completion, the user can:
\begin{itemize}
\item temporarily stop execution;
\item analyze the current state;
\item resume simulation;
\item or advance in a controlled manner.
\end{itemize}

This mode of use is especially useful in learning models, in initial debugging and when reading the behavior of elements whose function is not yet completely clear to the user.

\subsection{Stop and configure}
\label{subsec:cap5_parar_configurar}

The stop and configuration commands also deserve attention. Stopping the simulation is not just ending a process, but stopping the evolution of the model in a controlled way. Configuring the simulation means adjusting execution conditions before starting the experiment.

Although the depth of these parameters depends on the model and the context, the user must remember that the execution is not completely rigid: it can be prepared, started, observed and stopped according to different needs.

\section{Observing model behavior during execution}
\label{sec:cap5_observando_comportamento_execucao}

Observation of execution must be conducted in a guided manner. If the user tries to simultaneously follow all the GUI signals, the entire reading will become fuzzy. The ideal is to observe the model in levels.

\subsection{First level: the global flow}
\label{subsec:cap5_primeiro_nivel_fluxo_global}

At the first level, the user must observe the behavior of the model globally. The most important questions are:
\begin{itemize}
\item did the simulation start correctly?
\item does the main flow seem coherent with the model structure?
\item does the system evolve plausibly?
\item does termination of execution make sense?
\end{itemize}

This is a macroscopic reading. She doesn't look for fine details, but overall consistency.

\subsection{Second level: specific elements}
\label{subsec:cap5_segundo_nivel_elementos_especificos}

After observing global behavior, the user can focus on specific elements of the model. This may include:
\begin{itemize}
\item a component whose function is not yet clear;
\item a point in the flow where the behavior seems particularly important;
\item a relevant queue, resource or supporting data;
\item a property whose effect has recently been modified.
\end{itemize}

This localized reading is more productive after the general dynamics have already been minimally understood.

\subsection{Third level: comparison between runs}
\label{subsec:cap5_terceiro_nivel_comparacao_execucoes}

An especially useful practice is to compare the model run before and after small changes. This procedure helps the user understand not only how the system works, but also how it will react to changes. The GUI, in this sense, is not just an observation interface, but also a comparison laboratory.

\section{The \textit{Simulang} Tab as a Textual Representation of the Model}
\label{sec:cap5_aba_simulang}

After the reader has already seen the model in the GUI and observed it in execution, the textual representation starts to make much more sense. This is exactly why the \textit{Simulang} tab appears now, and not before.

The function of the \textit{Simulang} tab is not to replace the GUI, but to offer a second way of reading the same model. In this way of reading, the user stops seeing only the graphical organization of the system and starts seeing a corresponding textual description. This is valuable for several reasons:

\begin{itemize}
\item makes the model more formally explicit;
\item allows you to compare the visual representation with the textual representation;
\item helps consolidate understanding of the model;
\item paves the way for later chapters on language and development.
\end{itemize}

\subsection{How to Read the \textit{Simulang} Tab}
\label{subsec:cap5_como_ler_simulang}

Reading the \textit{Simulang} tab must be done with the same didactic caution used in the rest of the first part of the book. The user does not need, at this point, to master the entire syntax or semantics of the language. What he needs is to accomplish that:
\begin{itemize}
\item the textual description corresponds to the model you have already seen in the GUI;
\item visual elements of the model have a textual expression;
\item the language offers a more formal view of the modeled system.
\end{itemize}

The most productive way to start this reading is to compare excerpts from the textual representation with already known elements of the graphical interface. Instead of trying to decipher the entire text at once, the user should ask:
\begin{quote}
What part of the graphical model does this region of the text seem to describe?
\end{quote}

\subsection{Why textual representation is useful to the user}
\label{subsec:cap5_por_que_simulang_util_ao_usuario}

Even if the user continues to work primarily through the GUI, the textual representation offers concrete benefits. It helps to:
\begin{itemize}
\item consolidate understanding of the model;
\item make the system structure more explicit;
\item understand the model beyond its visual layout;
\item start thinking about the relationship between interface and language.
\end{itemize}

This experience is particularly valuable because it shows, from an early stage, that \textit{GenESyS} does not treat the model just as a drawing, but also as a formal description.

\section{Tracing, breakpoints and initial observability}
\label{sec:cap5_tracing_breakpoints_observabilidade}

An important feature of a well-structured simulator is to offer mechanisms for observing the behavior of the running system. In \textit{GenESyS}, this is expressed through different means available in the GUI, including traces, events, reports, and controlled stopping or tracking mechanisms.

At this stage of the manual, the reader does not yet need to deeply master these mechanisms. But you need to understand their general function: they make the simulation more observable.

\subsection{Tracing as reading behavior}
\label{subsec:cap5_tracing_leitura_comportamento}

Tracing should be seen as a way of recording what happens in the simulation. It helps the user follow the model not only through the visual perception of the interface, but also through an operational narrative of what is happening.

This is especially useful when the model's behavior is not immediately evident from the graphics area alone. In such cases, tracing answer helps:
\begin{itemize}
\item what events are occurring;
\item in what order;
\item with what apparent effects on the state of the system.
\end{itemize}

\subsection{Breakpoints as a controlled observation tool}
\label{subsec:cap5_breakpoints_observacao_controlada}

Breakpoints, when used, allow you to interrupt the simulation at points of interest. For the user, this turns the execution into an examineable object. Instead of letting the model evolve completely without intervention, you can stop at relevant states and take a closer look at what is happening.

This tool is particularly useful when:
\begin{itemize}
\item if you want to study the behavior of a specific part of the model;
\item if you want to validate the effect of a modification;
\item if you want to better understand the dynamics of a specific element.
\end{itemize}

\section{First reading of the results}
\label{sec:cap5_primeira_leitura_resultados}

In addition to observing the execution in progress, the user needs to start developing the habit of interpreting the results produced by the simulator. Here again, the guidance must be progressive.

The first reading of results should not be statistically sophisticated or overly ambitious. She must answer simple questions:
\begin{itemize}
\item did the simulation behave as expected?
\item were there coherent signs of progression and termination?
\item Did simple changes to the model produce observable effects?
\item Do the reports or signals presented by the GUI seem compatible with the model structure?
\end{itemize}

Only after this initial stage does it become useful to move on to more technical readings.

\subsection{Results as continuation of model reading}
\label{subsec:cap5_resultados_continuacao_leitura}

An important point is that the results should not be treated as something separate from the model. They are a continuation of reading the model by other means. The graphical model shows structure, the execution shows dynamics, and the results show observable consequences of that dynamic. These three dimensions must always be articulated.

\section{Best practices for observing running models}
\label{sec:cap5_boas_praticas_observacao_execucao}

Some practices make observing the much simulation more beneficial.

\subsection{Set an observation focus}
\label{subsec:cap5_defina_foco_observacao}

Before starting execution, try to decide what to observe first: the global flow, a specific component, an effect of a certain change, or the relationship between GUI and Simulang. This prevents dispersion.

\subsection{Switch between views}
\label{subsec:cap5_alterne_entre_visoes}

The strength of the \textit{GenESyS} GUI is in offering multiple perspectives on the same model. Switch between:
\begin{itemize}
\item graphic structure;
\item observed execution;
\item traces or events;
\item textual representation in \textit{Simulang}.
\end{itemize}

This alternation often produces deeper understanding than any single view.

\subsection{Make small changes and compare}
\label{subsec:cap5_pequenas_modificacoes_compare}

Whenever possible, make small, controlled changes to the model and compare the behavior before and after. This strategy turns observation into a guided experiment.

\subsection{Don't try to exhaust everything in a single run}
\label{subsec:cap5_nao_esgotar_tudo}

A simulation rarely reveals everything the model can teach in a single run. It is better to read successively, each with a clearer focus, than to try to absorb everything at once.

\begin{table}[htbp]
\centering
\caption{Recommended Strategy for Observing Running Models in GenESyS.}
 \label{tab:cap5_estrategia_observacao_execucao}
\small
 \begin{tabular}{|p{4.3cm}|p{6.9cm}|}
\hline
\textbf{Etapa} & \textbf{Objetivo} \\
\hline
Observe the graphical structure & Understand the static organization of the model \\
\hline
Start simulation & Transform structure into dynamic behavior \\
\hline
Monitor global flow & Check overall consistency of execution \\
\hline
Explore tracing and events & Make behavior more observable \\
\hline
Read the tab \textit{Simulang} & Relate the graphic model to its textual description \\
\hline
Compare runs & Understand the effect of small model changes \\
\hline
 \end{tabular}
\end{table}

\section{Final Considerations}
\label{sec:cap5_consideracoes_finais}

This chapter deepened the use of the \textit{GenESyS} GUI in an essential direction: the articulation between the model's visual structure, simulation execution, behavioral observability and textual representation in \textit{Simulang}. The main idea that the reader should remember is that these dimensions do not compete with each other. They complement each other. The graphical model allows you to see the organization of the system, the simulation allows you to observe its dynamics, the tracing and pausing mechanisms make this behavior more readable, and the \textit{Simulang} tab offers a textual way to consolidate this understanding.

With this, the first part of the manual begins to gain real density of use. The reader no longer just opens the GUI and loads a model; he began to observe the simulator in operation and report its various representations. The natural next step from here is to consolidate the use of the system with more guided examples and prepare the transition to a more complete understanding of the relationship between modeling, simulation and the internal structure of GenESyS.

Test your understanding of this content by answering the following questions:

\begin{enumerate}
\item Why should the execution of a model be understood as a transition from the static structure to the system dynamics?

\item What is the use of commands such as pause, resume and step-by-step during the simulation?

\item Why should the \textit{Simulang} tab only be introduced after the user has already seen and executed the model in the GUI?

\item How do tracing, breakpoints and results complement visual observation of model execution?
\end{enumerate}
